\documentclass[12pt]{article}

\usepackage[spanish]{babel}
\usepackage[utf8]{inputenc}
\usepackage[T1]{fontenc}
\usepackage{geometry}
\usepackage{array}
\usepackage{longtable}
\usepackage{hyperref}

\geometry{letterpaper,margin=2.5cm}

\begin{document}
	
	
	% Encabezado institucional
	\begin{center}
		{\bfseries UNIVERSIDAD AUTÓNOMA METROPOLITANA}\\[4pt]
		{\bfseries UNIDAD CUAJIMALPA}\\[4pt]
		DIVISIÓN DE CIENCIAS DE LA COMUNICACIÓN Y DISEÑO\\[4pt]
		{\bfseries LICENCIATURA EN TECNOLOGÍAS DE LA INFORMACIÓN}
	\end{center}
	
	\vspace{1em}
	
	\begin{table}[h]
		\renewcommand{\arraystretch}{1.2}
		\begin{tabular}{@{}l@{\hspace{0.5cm}}l@{}}
			PROFESOR: &
			\underline{\makebox[0.7\textwidth][c]{\MakeUppercase{Juan Manuel Alvarado López}}} \\[0.3em]
			ADSCRITO AL: &
			{\small \underline{\makebox[0.7\textwidth][l]{\textbf{DEPARTAMENTO DE TECNOLOGÍAS DE LA INFORMACIÓN}}}} \\[0.3em]
			OTRO: &
			\underline{\makebox[0.7\textwidth][c]{ }} \\[0.3em]
			UEA: &
			\underline{\makebox[0.7\textwidth][c]{\MakeUppercase{Inteligencia Artificial I}}} \\[0.3em]
		\end{tabular}
	\end{table}
	
	
	% Objetivo
	\noindent{OBJETIVO:} 
	
	Que al final del curso el alumno sea capaz de:
	\begin{enumerate}
		\item Explicar las características y objetivos, así como posibles aplicaciones de la inteligencia artificial.
		\item Emplear los métodos utilizados por la inteligencia artificial para desarrollar sistemas de cómputo.
	\end{enumerate}
	
	\vspace{1em}
	
	% HORARIO en tabla
	\section*{Horario}
	
	\begin{center}
		\renewcommand{\arraystretch}{1.3}
		\begin{tabular}{|c|c|c|}
			\hline
			\textbf{Día} & \textbf{Horario} & \textbf{Aula} \\
			\hline
			\underline{\makebox[2.8cm][c]{Martes}} &
			\underline{\makebox[2.8cm][c]{8--10}} &
			\underline{\makebox[2.8cm][c]{A-505}} \\
			\hline
			\underline{\makebox[2.8cm][c]{Jueves}} &
			\underline{\makebox[2.8cm][c]{8--11}} &
			\underline{\makebox[2.8cm][c]{A-505}} \\
			\hline
			%			\underline{\makebox[2.8cm][c]{\ }} &
			%			\underline{\makebox[2.8cm][c]{\ }} &
			%			\underline{\makebox[2.8cm][c]{\ }} \\
			%			\hline
		\end{tabular}
	\end{center}
	
	
	\vspace{1em}


% Plan de 11 semanas 
\section*{Plan de trabajo (11 semanas)}

\setlength{\extrarowheight}{2pt}
\begin{longtable}{|>{\centering\arraybackslash}p{1.2cm}|p{7.2cm}|p{7.2cm}|}
	\hline
	\textbf{Sem} & \textbf{Temas principales} & \textbf{Actividades (resumen)} \\
	\hline
	\endfirsthead
	\hline
	\textbf{Sem} & \textbf{Temas principales} & \textbf{Actividades (resumen)} \\
	\hline
	\endhead
	
	1 &
	Presentación general de la asignatura y explicación del programa operativo y plan de trabajo. &
	 \\
	\hline
	
	2 &
	Definición, historia y enfoques de la inteligencia artificial; introducción a los agentes inteligentes, su percepción, acción y racionalidad. &
	Debatir definiciones y casos actuales de IA; identificar agentes, entornos y elementos del modelo PEAS en situaciones reales. \\
	\hline
	
	3 &
	Propiedades de los entornos de tarea y tipos de agentes: reactivos simples, basados en modelo, metas, utilidad y aprendizaje. &
	Clasificar entornos de distintos problemas y programar un agente reactivo. \\
	\hline
	
	4 &
	Formulación de problemas de búsqueda: espacio de estados, nodos, funciones sucesoras y costos; búsqueda no informada mediante BFS, DFS y costo uniforme. &
	Modelar el 8-puzzle y laberintos como espacios de estados; implementar BFS en clase y DFS de forma individual. \\
	\hline
	
	5 &
	Búsquedas informadas, funciones heurísticas, búsqueda greedy best-first y algoritmo A*; admisibilidad y consistencia de heurísticas. &
	Diseñar y evaluar heurísticas. Implementar búsqueda greedy o algoritmo A*. \\
	\hline
	
	6 &
	Examen. &
	Revisión de examen. \\
	\hline
	
	7 &
	Métodos de búsqueda para optimización: Hill Climbing y Recocido Simulado. &
	Revisar ejemplos de optimización con ambos algoritmos y asignar el proyecto del curso con Hill Climbing o Recocido Simulado. \\
	\hline
	
	8 &
	Fundamentos de representación del conocimiento, requisitos y esquemas; redes semánticas &
	Representar un dominio pequeño con distintos esquemas y modelar conocimiento mediante frames. \\
	\hline
	
	9 &
	Sistemas basados en reglas de producción, inferencia hacia adelante y hacia atrás, ontologías, sistemas expertos y visión por computadora. &
	Desarrollar un mini sistema basado en reglas y revisar ejemplos de sistemas expertos y aplicaciones de visión por computadora. \\
	\hline
	
	10 &
	Introducción al procesamiento de lenguaje natural, robótica y machine learning; implicaciones éticas y sociales de la inteligencia artificial. &
	Analizar ejemplos de aplicaciones de NLP, robótica y machine learning; discutir casos relacionados con dilemas éticos y sociales. \\
	\hline
	
	11 &
	Impacto de la inteligencia artificial en individuos, organizaciones y sociedad: empleo, privacidad, regulación y aspectos legales. &
	Analizar y discutir ejemplos de impacto social y legal. Entrega del proyecto. \\
	\hline
	
	11bis &
	Examen final. &
	 \\
	\hline
	
\end{longtable}



	
	% Criterios de evaluación (desde el .md)
	\section*{Criterios e instrumentos de evaluación}
	
	La calificación final se integra a partir de los siguientes instrumentos de evaluación. Cada uno mide distintas dimensiones del aprendizaje: la comprensión conceptual, la aplicación práctica y el desempeño individual.
	
	\begin{center}
		\begin{tabular}{|l|p{7cm}|c|}
			\hline
			\textbf{Instrumento} & \textbf{Descripción} & \textbf{Porcentaje} \\
			\hline
			Proyecto & Trabajo individual o colaborativo entregado fuera del horario de clase & 25\% \\
			\hline
			Actividades & Ejercicios, tareas, prácticas, participación & 25\% \\
			\hline
			Examen parcial & Evaluación escrita y/o práctica a mitad del periodo & 25\% \\
			\hline
			Examen final & Evaluación escrita y/o práctica al cierre del periodo & 25\% \\
			\hline
			Total & & 100\% \\
			\hline
		\end{tabular}
	\end{center}
	
	\vspace{0.5em}
	
	\noindent\textbf{Calificación mínima aprobatoria.} La calificación mínima aprobatoria es \textbf{6.0 (seis)} en escala de 0 a 10.
	
	\noindent\textbf{Equivalencias numéricas:}
	\begin{itemize}
		\item NA: 0.0 a < 6.0
		\item S: 6.0 a < 8.0
		\item B: 8.0 a < 9.0
		\item MB: 9.0 a 10.0
	\end{itemize}
	
	\subsection*{Entrega de trabajos y actividades}
	
	\begin{itemize}
		\item Las fechas de entrega se publicarán con anticipación en la plataforma institucional y no serán modificadas salvo causa de fuerza mayor debidamente justificada.
		\item Los trabajos se entregarán \textbf{exclusivamente en los formatos indicados} (PDF, Word, ZIP u otro especificado) y a través del canal oficial del curso (por ejemplo, la plataforma LMS institucional).
		\item \textbf{No se aceptarán entregas mediante ligas o vínculos a servicios de almacenamiento en la nube} (Google Drive, OneDrive, Dropbox, etcétera).
		\item Los trabajos entregados después de la fecha y hora límite no serán considerados para evaluación, salvo circunstancias excepcionales documentadas.
	\end{itemize}
	
	% Política de asistencia (desde el .md)
	\section*{Política de asistencia}
	
	La asistencia regular es un factor determinante en el aprendizaje. El seguimiento de las actividades presenciales, la participación en clase y la interacción entre pares no pueden reemplazarse por otros medios.
	
	\subsection*{Requisito mínimo de asistencia}
	
	\textbf{Se espera una asistencia regular a las sesiones del curso.} Cuando la asistencia mínima establecida no se cumple, el estudiante deja de ser elegible para consideraciones especiales en entregas extemporáneas, reprogramación de evaluaciones u otros ajustes extraordinarios.
	
	\subsection*{Definición de asistencia, retardo y falta}
	
	\begin{center}
		\begin{tabular}{|l|p{8cm}|}
			\hline
			\textbf{Condición} & \textbf{Criterio} \\
			\hline
			Asistencia & Presente dentro de los primeros 15 minutos de iniciada la sesión \\
			\hline
			Retardo & Presente entre el minuto 16 y el minuto 30 de iniciada la sesión \\
			\hline
			Falta & Presente después del minuto 30, o no se presenta \\
			\hline
			Falta por abandono & Salida del aula sin justificación por más de 20 minutos continuos durante la sesión \\
			\hline
		\end{tabular}
	\end{center}
	
	\vspace{0.5em}
	
	\noindent 2 retardos equivalen a 1 falta. Se permiten hasta \textbf{3 faltas}. A partir de la cuarta falta, el estudiante pierde el derecho a consideraciones especiales.
	
	\subsection*{Faltas justificadas}
	
	Las inasistencias por motivos de salud, compromisos institucionales (representación deportiva, cultural o académica) o causas de fuerza mayor deberán notificarse \textbf{antes de la sesión o a más tardar 48 horas después}, acompañadas del documento que las acredite. \textbf{Sin documento de evidencia no se justifica la falta.}
	
	% Integridad académica (desde el .md)
	\section*{Integridad académica: plagio y uso de IA}
	
	La formación en este curso busca desarrollar competencias auténticas. La elaboración propia de los trabajos y evaluaciones es parte fundamental de ese proceso.
	
	\subsection*{Plagio}
	
	Se considera plagio toda reproducción total o parcial de ideas, textos, código fuente, diagramas u otros contenidos de terceros, sin el reconocimiento explícito de la fuente original mediante cita adecuada.
	
	\begin{itemize}
		\item El estudiante es responsable de citar correctamente cualquier fuente consultada, ya sea impresa, digital o generada por herramientas de terceros.
		\item Los trabajos que presenten indicios de plagio serán reportados y \textbf{no serán tomados en cuenta para efectos de evaluación}, quedando con valor de cero en el instrumento correspondiente.
	\end{itemize}
	
	\subsection*{Uso de herramientas de inteligencia artificial}
	
	El uso de herramientas de inteligencia artificial generativa (como ChatGPT, Copilot, Gemini u otras similares) está \textbf{sujeto a las indicaciones específicas de cada actividad}. En términos generales:
	
	\begin{itemize}
		\item \textbf{Uso permitido y declarado:} Las herramientas de IA podrán utilizarse como apoyo para exploración, estructuración o revisión de ideas, siempre que el estudiante declare explícitamente su uso al momento de la entrega.
		\item \textbf{Uso no permitido:} Queda prohibido presentar como producción propia contenido íntegramente generado por una IA, sin intervención analítica o creativa sustancial del estudiante.
		\item \textbf{En exámenes:} Queda estrictamente prohibido el uso de cualquier herramienta de IA, navegador web u otro recurso externo no autorizado explícitamente por el docente.
	\end{itemize}
	
	\subsection*{Consideraciones generales de integridad académica}
	
	Se hará uso de las \textbf{herramientas de detección de plagio y contenido generado por IA} disponibles. Si se detecta plagio o uso indebido de IA en un trabajo o actividad evaluable, dicho instrumento \textbf{quedará sin efecto para efectos de calificación} y se registrará con valor de cero, sin posibilidad de reentrega. En exámenes, el uso indebido de dispositivos electrónicos, o de cualquier material no autorizado o solicitado, anula la calificación del examen y el caso será reportado a las instancias académicas correspondientes. 
	
	
	% Conducta y convivencia (desde el .md)
	\section*{Conducta y convivencia en el aula}
	
	Las siguientes normas de convivencia aplican durante todas las sesiones presenciales:
	
	\begin{itemize}
		\item \textbf{Puntualidad:} Se espera la presencia de los estudiantes al inicio de cada sesión con tolerancia de 15 minutos.
		\item \textbf{Respeto:} Tener trato respetuoso hacia los compañeros de clase y el personal de la institución en todo momento.
		\item \textbf{Dispositivos electrónicos:} El uso de computadoras, tabletas y teléfonos inteligentes está \textbf{permitido únicamente con fines académicos} relacionados con la sesión en curso.
		\item \textbf{Silencio y atención:} Se solicita silenciar notificaciones y evitar conversaciones paralelas que interrumpan la dinámica de la clase.
		\item \textbf{Cuidado de instalaciones:} Los estudiantes son responsable de mantener en buen estado el mobiliario, equipo y espacios del aula.
	\end{itemize}
	
	
	
\end{document}